\section{Introduction}

\lettrine[lines=2]{\color{titleblue}L}{ es} réseaux de distribution d'eau constituent des infrastructures critiques pour les environnements urbains, en assurant une alimentation continue en eau potable pour les usagers résidentiels, industriels et commerciaux. Cependant, ces systèmes sont fortement exposés aux fuites en raison du vieillissement des infrastructures, des fluctuations de pression et de facteurs environnementaux. Les fuites non détectées peuvent entraîner des pertes d'eau importantes, une hausse des coûts d'exploitation, une baisse de l'efficacité du service et des dommages structurels à long terme. Par conséquent, le développement de systèmes fiables et automatisés de détection des fuites est devenu un enjeu majeur de recherche et d'industrie.

Les méthodes traditionnelles de détection des fuites reposent généralement sur l'inspection manuelle, la détection acoustique, ou des techniques de seuillage basées sur des règles appliquées aux mesures de pression et de débit. Bien que ces approches puissent être efficaces dans des scénarios contrôlés, elles manquent souvent de passage à l'échelle et d'adaptabilité dans des réseaux urbains complexes. De plus, ces méthodes peinent à capturer les dynamiques temporelles et les motifs non linéaires associés aux événements de fuite, en particulier en présence de bruit et de conditions d'exploitation variables.

Avec le déploiement croissant des réseaux de capteurs et des technologies de l'Internet des objets (IoT) dans les systèmes d'eau intelligents, de grands volumes de séries temporelles multivariées sont désormais disponibles. Cela ouvre la voie à des approches orientées données, en particulier les méthodes d'apprentissage profond, qui ont démontré de fortes capacités pour la modélisation séquentielle. Des architectures telles que les réseaux de neurones convolutionnels (CNN) et les réseaux Long Short-Term Memory (LSTM) sont largement utilisées pour la classification de séries temporelles grâce à leur capacité à capturer respectivement les dépendances locales et de long terme. Toutefois, les approches existantes s'appuient souvent sur des architectures monomodèles et manquent d'évaluations complètes face à des méthodes alternatives, ce qui limite leur capacité à exploiter pleinement les complémentarités de modélisation.

Dans ce travail, nous proposons une architecture hybride d'apprentissage profond combinant CNN, LSTM bidirectionnel (BiLSTM) et mécanisme d'auto-attention pour la détection automatique des fuites d'eau. Le composant CNN extrait des caractéristiques temporelles locales à partir des signaux capteurs, tandis que le BiLSTM capture les dépendances bidirectionnelles dans le temps. L'intégration d'une couche d'auto-attention permet au modèle de se concentrer sur les segments les plus informatifs de la séquence d'entrée, améliorant à la fois les performances de détection et l'interprétabilité.

Au-delà de l'architecture du modèle, nous concevons une chaîne de traitement complète adaptée à un déploiement réel. Celle-ci inclut une ingénierie de caractéristiques guidée par le domaine, une construction de séquences par fenêtre glissante, et une séparation temporelle stricte entre les ensembles d'entraînement, de calibration et de test afin d'éviter les fuites de données. En outre, nous traitons le déséquilibre de classes inhérent à la détection des fuites en adoptant des métriques comme l'AUC-PR et en effectuant une calibration du seuil pour contrôler le compromis entre sensibilité de détection et taux de fausses alertes.

Afin de garantir une évaluation rigoureuse, le modèle proposé est comparé à plusieurs approches de référence en machine learning et en apprentissage profond. Cette analyse de benchmark fournit une évaluation complète de l'efficacité de l'architecture hybride par rapport aux méthodes alternatives.

Les résultats expérimentaux obtenus sur un jeu de données réel montrent que l'approche proposée atteint de solides performances, avec des scores AUC-ROC et AUC-PR élevés, ainsi qu'un compromis favorable entre précision et rappel. En particulier, le modèle maintient un faible taux de fausses alertes tout en conservant une forte capacité de détection, ce qui le rend adapté aux systèmes de surveillance pratiques où fiabilité et réactivité sont essentielles.

Les principales contributions de cet article peuvent être résumées comme suit :
\begin{itemize}
    \item Une architecture hybride CNN-BiLSTM renforcée par un mécanisme d'auto-attention pour la détection des fuites dans des séries temporelles multivariées.
    \item Une chaîne robuste de traitement des données incluant l'ingénierie de caractéristiques, la transformation par fenêtre glissante et un découpage temporel sans fuite de données.
    \item Un cadre d'évaluation prenant en compte le déséquilibre de classes, intégrant l'AUC-PR et la calibration du seuil pour la prise de décision opérationnelle.
    \item Une comparaison de benchmark complète face à des modèles classiques et de deep learning afin de valider l'efficacité de l'approche proposée.
\end{itemize}